\documentclass[11pt, a4paper]{article}

% ── Packages ──────────────────────────────────────────────────────────
\usepackage[utf8]{inputenc}
\usepackage[T1]{fontenc}
\usepackage{lmodern}
\usepackage[margin=1in]{geometry}
\usepackage{amsmath, amssymb, amsthm}
\usepackage{booktabs}
\usepackage{longtable}
\usepackage{graphicx}
\usepackage{hyperref}
\usepackage{xcolor}
\usepackage{listings}
\usepackage{tcolorbox}
\usepackage{polya}

\hypersetup{
  colorlinks=true,
  linkcolor=polyablue,
  urlcolor=polyablue,
  citecolor=polyablue,
}

% ── Code listing style ────────────────────────────────────────────────
\lstdefinestyle{polya-python}{
  language=Python,
  basicstyle=\ttfamily\small,
  keywordstyle=\color{polyablue}\bfseries,
  commentstyle=\color{polyagray}\itshape,
  stringstyle=\color{polyagreen},
  numbers=left,
  numberstyle=\tiny\color{polyagray},
  numbersep=8pt,
  frame=single,
  framerule=0.4pt,
  rulecolor=\color{polyagray},
  breaklines=true,
  showstringspaces=false,
  tabsize=4,
}
\lstset{style=polya-python}
\title{Break-Even Analysis}
\author{uber-polya}
\date{2026-02-26}

\begin{document}

\maketitle
\tableofcontents
\newpage

% ── Phase A: Problem Formulation ─────────────────────────────────────
\section{Problem Formulation}

\begin{polyamodel}

\textbf{Problem Type}: Find \hfill
\textbf{Domain}: Numerical Methods


\end{polyamodel}

% ── Universe ──────────────────────────────────────────────────────────
\subsection{Universe}

\begin{itemize}
  \item $See problem description$
\end{itemize}

% ── Variables ─────────────────────────────────────────────────────────
\subsection{Variables}

\begin{longtable}{lll}
\toprule
\textbf{Symbol} & \textbf{Meaning} & \textbf{Type / Range} \\
\midrule
\endhead
$x$ & decision variable & $\mathbb{{R}}$ \\
\bottomrule
\end{longtable}

% ── Structure ─────────────────────────────────────────────────────────
\subsection{Structure}

A small business is considering launching a new product. Fixed costs are \$25,000 (tooling, marketing). Variable cost per unit is \$12. Selling price per unit is \$35. What is the break-even quantity? How does the break-even point shift if the price drops 10\% or variable costs increase 20\%?

% ── Mapping ───────────────────────────────────────────────────────────
\subsection{Real-World Mapping}

\begin{longtable}{ll}
\toprule
\textbf{Real-World Concept} & \textbf{Mathematical Object} \\
\midrule
\endhead
problem input & $instance parameters$ \\
\bottomrule
\end{longtable}

% ── Constraints ───────────────────────────────────────────────────────
\subsection{Constraints}

\begin{enumerate}
  \item $Break-even analysis$
  \hfill \textit{(finding the production volume where profit turns positive)}
  \item $Symbolic computation$
  \hfill \textit{(deriving exact formulas with SymPy)}
  \item $Sensitivity analysis$
  \hfill \textit{(understanding how parameter changes affect the result)}
  \item $Contribution margin$
  \hfill \textit{(Price minus Variable cost per unit)}
  \item $Business decision modeling$
  \hfill \textit{(translating financial questions into mathematical models)}
\end{enumerate}

% ── Objective / Claim ─────────────────────────────────────────────────
\subsection{Claim}


% ── Approach ──────────────────────────────────────────────────────────
\subsection{Approach}

\begin{description}
  \item[Suggested approach:] Symbolic (SymPy) + Numerical (NumPy) Break-Even Analysis
  \item[Complexity class:] 
  \item[Available tools:] Symbolic
\end{description}
% ── Phase B: Solution ────────────────────────────────────────────────
\section{Solution}

\begin{polyaresult}

\textbf{Answer}: Solution computed


\textbf{Optimal}: No / Unknown \hfill
\textbf{Feasible}: Yes
\textbf{Algorithm}: Symbolic (SymPy) + Numerical (NumPy) Break-Even Analysis \quad $$

\textbf{Time}: 0.11466744501376525s


\end{polyaresult}

% ── Solution Details ──────────────────────────────────────────────────
\subsection{Solution Details}


Method: Symbolic algebra via SymPy to derive the closed-form break-even formula, then numerical evaluation with NumPy for sensitivity analysis. Break-even occurs where Revenue(Q) = Cost(Q), i.e., P*Q = F + V*Q, yielding Q* = F / (P - V). Sensitivity sweeps over price and variable cost perturbations.

% ── Verification ──────────────────────────────────────────────────────
\subsection{Verification}

\begin{longtable}{lcl}
\toprule
\textbf{Check} & \textbf{Status} & \textbf{Detail} \\
\midrule
\endhead
Feasibility & \passmark & All constraints satisfied \\
\bottomrule
\end{longtable}

% ── Phase C: Interpretation ──────────────────────────────────────────
\section{Interpretation}

% ── The Question & The Answer ─────────────────────────────────────────
\subsection{The Question}
A small business is considering launching a new product.

\subsection{The Answer}

\begin{polyainsight}[title={Bottom Line}]
A feasible solution was found.
\end{polyainsight}

% ── What This Means ───────────────────────────────────────────────────
\subsection{What This Means}

- Exact symbolic break-even quantity: Q = Fixed / (Price - Variable)
- Base case: \textasciitilde{}1087 units
- Sensitivity table showing break-even under price drops and cost increases
- Profit projections at various production volumes

Symbolic algebra via SymPy to derive the closed-form break-even formula, then numerical evaluation with NumPy for sensitivity analysis. Break-even occurs where Revenue(Q) = Cost(Q), i.e., P*Q = F + V*Q, yielding Q* = F / (P - V). Sensitivity sweeps over price and variable cost perturbations.

% ── Sensitivity Analysis ──────────────────────────────────────────────
\subsection{Sensitivity Analysis}

\begin{longtable}{llrrl}
\toprule
\textbf{Parameter} & \textbf{Current} & \textbf{Change} & \textbf{New Obj.} & \textbf{Class} \\
\midrule
\endhead
worst\_case & present & removed &  &
\textcolor{polyagreen}{robust} \\
\bottomrule
\end{longtable}

% ── Figures ───────────────────────────────────────────────────────────

% ── Recommendations ───────────────────────────────────────────────────
\subsection{Recommendations}

\begin{enumerate}
  \item Review the model assumptions to ensure real-world fidelity.
\end{enumerate}

% ── Limitations ───────────────────────────────────────────────────────
\subsection{Limitations}

\begin{polyawarnbox}
\begin{itemize}
  \item Model assumes deterministic parameters; real-world variability not captured.
  \item Solution quality depends on the accuracy of input data.
\end{itemize}
\end{polyawarnbox}

% ── Appendix ─────────────────────────────────────────────────────────

\end{document}